
\subsection{3D Scene Synthesis}
\label{subsec:p3_s8_2_2}


Scene-level synthesis extends 3D generation from a single asset $\mathbf{r}\!\in\!\mathcal{R}$ to an entire environment: a clutter of objects whose positions, orientations, and pairwise relations together compose a navigable room or open-world scene.
Where §\ref{subsec:p3_s8_2_1} models the conditional distribution of one object given a prompt, scene synthesis must jointly model (i) a variable number of objects, (ii) their global $\mathrm{SE}(3)$ placements, (iii) their pairwise spatial and functional relations, and (iv) the envelope geometry (floor, walls, ceiling) that closes the scene; the output is a structured set rather than a vector in $\mathcal{R}$, and the conditioning broadens from text to layouts, scene graphs, sparse calibrated views, and full camera trajectories.
The result must be navigable, not only renderable, which makes the problem harder than object synthesis along three axes at once: geometric coherence across the clutter, semantic consistency between related objects, and view generalization along an $\mathrm{SE}(3)$ trajectory rather than a fixed look-at orbit.
Four families answer it: \emph{layout- and scene-graph-conditioned} generators that produce the structured set from a symbolic prior; \emph{view-to-3D} synthesizers that regress or sample a continuous radiance or Gaussian field from sparse calibrated views; \emph{view-to-view} synthesizers that skip the explicit 3D representation and learn a conditional map from observed to novel views; and \emph{video-model-based} synthesizers that treat a pretrained video diffusion model as an implicit 3D prior with the camera trajectory as an additional conditioning channel.


\subsubsection{Task Definition}
\label{subsubsec:p3_s8_2_2_task}

We cast 3D scene synthesis as a scene-level extension of the object-synthesis framework of §\ref{subsec:p3_s8_2_1}.
Let $\mathcal{S}$ denote the scene-output space whose elements
$\mathbf{s} = \big(\{\mathbf{r}_i\}_{i=1}^N,\, \{\mathbf{T}_i\}_{i=1}^N,\, \mathbf{G},\, \mathcal{E}_{\text{room}}\big)$
package $N$ object representations $\mathbf{r}_i\!\in\!\mathcal{R}$ (with $\mathcal{R}$ the per-object space of §\ref{subsec:p3_s8_2_1}), their global poses $\mathbf{T}_i\!\in\!\mathrm{SE}(3)$, a relational graph $\mathbf{G}$ over the object set, and an envelope geometry $\mathcal{E}_{\text{room}}$ encoding the room shell (walls, floor, ceiling).
The conditioning space broadens to
$\mathcal{C}^{\text{scene}} \supseteq \{\,\text{text }\mathbf{x},\,\text{layout }\mathbf{L},\,\text{scene graph }\mathbf{G}_{\text{in}},\,\text{sparse views }\{(I_v,\pi_v)\}_v,\,\text{camera trajectory }\{\pi_v\}_v\,\}$,
with a condition encoder $\kappa^{\text{scene}}:\mathcal{C}^{\text{scene}}\!\to\!\mathcal{C}^{\text{scene}\prime}$ mapping each input mode into a shared embedding.
The core learning problem is to train a conditional scene generator
\begin{equation}
G_\theta^{\text{scene}}:\mathcal{C}^{\text{scene}\prime}\rightarrow \mathcal{S},
\qquad
\hat{\mathbf{s}}\sim p_\theta^{\text{scene}}\!\big(\mathbf{s}\,\big|\,\kappa^{\text{scene}}(\mathbf{c})\big),
\label{eq:p3_s8_2_2_objective}
\end{equation}
so that the sampled scene $\hat{\mathbf{s}}\!\in\!\mathcal{S}$ admits coherent inter-object geometry, semantically plausible relations, and view-consistent renderings under arbitrary $\mathrm{SE}(3)$ trajectories.
Two limiting cases fix the scope of Eq.~\eqref{eq:p3_s8_2_2_objective}.
When $N\!=\!1$, $\mathbf{G}\!=\!\varnothing$, and $\mathcal{E}_{\text{room}}\!=\!\varnothing$, the scene degenerates to a single asset and the sampler recovers Eq.~\eqref{eq:p3_s8_2_1_objective}: scene synthesis is a strict generalization of object synthesis.
At the opposite extreme, when $\mathcal{S}$ is instantiated as an ordered sequence of views $\{I_v\}_{v=1}^V$ \emph{without an explicit 3D substrate}, the sampler reduces to a camera-trajectory-conditioned video generator (the regime of §\ref{para:p3_s8_2_2_video} below), so that ``scene synthesis'' and ``consistent video generation'' are two boundary regimes of the same problem.


\subsubsection{Methodology}
\label{subsubsec:p3_s8_2_2_methodology}

Each of the four paradigms is a different factorization of $G_\theta^{\text{scene}}$ in Eq.~\eqref{eq:p3_s8_2_2_objective}, from sparse structural controls (layouts and scene graphs) through continuous 3D representations to video-based synthesis; Figure~\ref{fig:p3_s8_2_2_paradigms} previews them against a common scene.

\begin{figure}[!ht]
    \centering
    \includegraphics[width=\linewidth]{figures/part3/section8/Four_paradigm_families_for_3D_scene_synthesis.pdf}
    \caption{Four paradigm families for 3D scene synthesis (Eq.~\eqref{eq:p3_s8_2_2_objective}): (a) layout- or scene-graph-conditioned placement of objects, (b) view-to-3D reconstruction into a radiance field or 3D Gaussians, (c) view-to-view synthesis with no explicit 3D, and (d) video-model-based synthesis along a camera path. The families are not mutually exclusive, and view consistency does not guarantee a valid 3D asset; room images are illustrative.}
    \label{fig:p3_s8_2_2_paradigms}
\end{figure}

\paragraph{Layout and Scene-graph Conditioned Methods.}
\label{para:p3_s8_2_2_layout}
A major line of research generates 3D environments from high-level structural priors such as layouts $\mathbf{L}$ or scene graphs $\mathbf{G}_{\text{in}}$, instantiating Eq.~\eqref{eq:p3_s8_2_2_objective} by materializing the structured tuple as an explicit object set,
\begin{equation}
\hat{\mathbf{s}} = \big\{(\hat{\mathbf{r}}_i,\,\hat{\mathbf{T}}_i)\big\}_{i=1}^N,
\qquad
\hat{\mathbf{s}}\sim p_\theta\!\Big(\{(\mathbf{r}_i,\mathbf{T}_i)\}_i \,\Big|\, \kappa(\mathbf{L} \text{ or } \mathbf{G}_{\text{in}}),\,\phi(\mathbf{x})\Big),
\label{eq:p3_s8_2_2_layout_objective}
\end{equation}
where the joint over the $N$ slots and their poses is factorized autoregressively (a transformer over object slots) or as a denoising diffusion over the slot tensor.
SceneFormer~\cite{wang2021sceneformer} explored autoregressive modeling of room layouts and furniture placement, and DiffuScene~\cite{tang2024diffuscene} introduced diffusion over object arrangements, improving realism and diversity.
Recent work adds language and visual reasoning for compositional control: LayoutGPT~\cite{feng2023layoutgpt} plans the layout with an LLM before neural scene construction, and LayoutVLM~\cite{sun2025layoutvlm} optimizes layout consistency under multimodal alignment.
Scene-graph approaches (GraphDreamer~\cite{gao2024graphdreamer}, CommonScenes~\cite{zhai2023commonscenes}, EchoScene~\cite{zhai2024echoscene}) use relational reasoning to keep placements plausible, and InstructScene~\cite{lin2024instructscene} adds natural-language conditions for controllable scene creation.
Beyond indoor settings, CityCraft~\cite{deng2024citycraft} scales layout synthesis to the city level, and GALA3D~\cite{zhou2024gala3d} couples layout-guided diffusion with Gaussian-based rendering for photorealistic multi-object environments.

\paragraph{View-to-3D Representation Synthesis.}
\label{para:p3_s8_2_2_view3d}
The second paradigm trades the explicit structured set for a \emph{continuous} 3D representation, recovering the scene as a radiance field or 3D Gaussian collection from sparse calibrated views; the scene tuple collapses into a single field,
\begin{equation}
\hat{\mathcal{R}}^{\text{scene}} = \mathcal{F}_\psi\!\Big(\big\{(I_v,\pi_v)\big\}_{v=1}^V\Big),
\qquad \hat{\mathbf{s}}\equiv \hat{\mathcal{R}}^{\text{scene}},
\label{eq:p3_s8_2_2_view3d_objective}
\end{equation}
where $\mathcal{F}_\psi$ is either a deterministic feed-forward regressor or a generative sampler over the field.
This scales the multi-view-and-lift recipe of §\ref{para:p3_s8_2_1_mv} from one asset to a full scene by replacing the per-object reconstructor with a scene-scale one.
Deterministic regressors such as PixelNeRF~\cite{yu2021pixelnerf}, PixelSplat~\cite{charatan2024pixelsplat}, DepthSplat~\cite{xu2024depthsplat}, and MVSplat~\cite{chen2024mvsplat} / MVSplat360~\cite{chen2024mvsplat360} stay strictly faithful to the input views and are reconstruction methods, treated in §\ref{para:p3_s8_2_3_feedforward}.
The generative variant samples the field instead: DreamScene~\cite{li2024dreamscene} introduces a diffusion prior over Gaussian fields for text-conditioned scene creation, and Generative 3DGS~\cite{Schwarz2025ggs} generates 3D Gaussians for scenes with video priors.

\paragraph{View-to-View Synthesis (Regression and Generative).}
\label{para:p3_s8_2_2_viewview}
The third paradigm skips the explicit 3D representation altogether and formulates novel view synthesis as a direct conditional 2D-to-2D map from observed to target views,
\begin{equation}
\hat{I}_{\pi^*}\sim p_\theta\!\Big(I_{\pi^*}\,\Big|\,\big\{(I_v,\pi_v)\big\}_{v=1}^V,\,\pi^*\Big),
\label{eq:p3_s8_2_2_viewview_objective}
\end{equation}
where $\pi^*$ is the target camera pose and no intermediate 3D representation is materialized.
Regression-based approaches learn a deterministic map: SRT~\cite{sajjadi2022scene} encodes the scene into latent tokens that support interpolation and extrapolation across viewpoints, and LVSM~\cite{jin2025lvsm} predicts novel views with a large feed-forward transformer and minimal 3D inductive bias.
Generative approaches model the conditional distribution over unseen views, which permits stochastic synthesis under ambiguity and large viewpoint changes: Photoconsistent-NVS~\cite{yu2023long} and SEVA~\cite{zhou2025stable} use diffusion to improve diversity and robustness along extended camera trajectories, and CausNVS~\cite{kong2025causnvs} generates views autoregressively and causally, modeling inter-view dependencies explicitly to stay coherent across complex viewpoint transitions.

\paragraph{Video-model Based Scene Synthesis.}
\label{para:p3_s8_2_2_video}
The fourth paradigm adds \emph{temporal} coherence to the viewpoint variety of the view-to-view family by exploiting large video generation models as spatiotemporal priors; the scene tuple is materialized as a frame sequence,
\begin{equation}
\{\hat{I}_v\}_{v=1}^V \sim p_\theta^{\text{video}}\!\Big(\{I_v\}_{v=1}^V \,\Big|\, \{\pi_v\}_{v=1}^V,\,\phi(\mathbf{x})\Big),
\label{eq:p3_s8_2_2_video_objective}
\end{equation}
treating the scene as the implicit 3D structure of a coherent video and the camera trajectory $\{\pi_v\}$ as an explicit conditioning channel.
Where the feed-forward reconstructors of §\ref{para:p3_s8_2_1_lrm} output a static asset in one pass, the video paradigm trades static 3D for temporally coherent multi-frame output that mirrors the navigability requirement of a scene more directly.
CameraCtrl~\cite{he2024cameractrl} and MotionCtrl~\cite{wang2024motionctrl} condition video generation on camera trajectories directly; ViewCrafter~\cite{yu2024viewcrafter} and Gen3C~\cite{ren2025gen3c} add geometric guidance by unprojecting observed views into an intermediate 3D representation and reprojecting it to the target viewpoints, which improves pose fidelity; and CAT4D~\cite{wu2025cat4d} and 4DiM~\cite{watson2024controlling} train jointly on posed static multi-view data and on temporal video data, so that viewpoint variation and temporal dynamics are modeled together for 4D scene synthesis.


\subsubsection{Benchmarks}
\label{subsubsec:p3_s8_2_2_benchmarks}

The four paradigms draw on different training corpora and report on different metric panels, each inherited from a precursor task.

\paragraph{Training Corpora.}
Table~\ref{tab:p3_s8_2_2_corpora} lists the corpora in use.
Layout-conditioned methods rely overwhelmingly on 3D-FRONT~\cite{fu20213dfront}, usually combined with its companion 3D-FUTURE~\cite{fu20213dfuture} furniture library for instance-level geometry.
View-based methods (P2 / P3) draw on real-world scans of residential interiors (ScanNet~\cite{dai2017scannet}, Matterport3D~\cite{chang2017matterport3d}), on photorealistic synthetic corpora with ground-truth depth and normals (Hypersim~\cite{roberts2021hypersim}), and on Replica~\cite{straub2019replica} for compact high-fidelity indoor benchmarks.
Video-model-based methods (P4) inherit the web-scale supervision of their underlying video prior, far larger than any 3D corpus but without explicit 3D supervision, supplemented by camera-pose-annotated subsets distilled from ScanNet, RealEstate10K~\cite{zhou2018stereo}, or hand-labeled trajectories.

\begin{table}[t]
    \centering
    \fontsize{8}{10}\selectfont
    \setlength{\tabcolsep}{1.0mm}
    \caption{Representative training corpora for 3D scene synthesis.}
    \label{tab:p3_s8_2_2_corpora}
    \begin{tabular}{L{0.20\textwidth} C{0.07\textwidth} C{0.14\textwidth} L{0.22\textwidth} L{0.29\textwidth}}
        \toprule
        \textbf{Dataset} & \textbf{Year} & \textbf{Scale} & \textbf{Acquisition} & \textbf{Primary use / paradigm} \\
        \midrule
        ScanNet~\cite{dai2017scannet} & 2017 & ${\sim}1.5\mathrm{K}$ scenes & Real-world RGB-D scans & Sparse-view supervision (P2 / P3) \\
        Matterport3D~\cite{chang2017matterport3d} & 2017 & ${\sim}90$ buildings & Real-world RGB-D scans & Whole-house multi-view (P2 / P3) \\
        Replica~\cite{straub2019replica} & 2019 & ${\sim}18$ scenes & Real-world high-fidelity scans & Photorealistic indoor benchmarks (P2 / P3) \\
        Hypersim~\cite{roberts2021hypersim} & 2021 & ${\sim}460$ scenes & Photorealistic synthetic & Ground-truth depth / normals (P2 / P3) \\
        3D-FRONT~\cite{fu20213dfront} & 2021 & ${\sim}19\mathrm{K}$ rooms & Synthetic indoor layouts & Layout / object-arrangement supervision (P1) \\
        3D-FUTURE~\cite{fu20213dfuture} & 2021 & ${\sim}10\mathrm{K}$ assets & Synthetic furniture library & Companion to 3D-FRONT (P1) \\
        \bottomrule
    \end{tabular}
\end{table}

\paragraph{Evaluation Benchmark Suites.}
Layout-conditioned methods report the metrics introduced with ATISS on 3D-FRONT (object-category frequency, pairwise-relation accuracy, and FID on rendered top-down or perspective views).
View-to-3D and view-to-view methods adopt the NeRF-era panel of PSNR / SSIM / LPIPS on held-out views of ScanNet, RealEstate10K~\cite{zhou2018stereo}, or DL3DV-10K~\cite{ling2024dl3dv}.
Video-model-based methods add FVD and a camera-trajectory-adherence metric, the mean translation and rotation error between requested and recovered camera poses.

\paragraph{Constraint-specific Metrics.}
Four axes mirror the scene tuple of Eq.~\eqref{eq:p3_s8_2_2_objective}: \textbf{layout adherence} (object-count agreement, pairwise-relation accuracy on $\mathbf{G}$, and bounding-box IoU between requested and synthesized placements); \textbf{geometric fidelity} where ground truth exists (Chamfer Distance, F-score, and depth and normal consistency against held-out views); \textbf{photometric quality} (PSNR / SSIM / LPIPS against ground-truth views, FID / KID against held-out scene distributions, and FVD for video outputs); and \textbf{camera-pose adherence}, most relevant for P4 (translation and rotation error of poses recovered by structure-from-motion or a learned pose regressor).
A paper reports the subset of axes its paradigm can be scored on, which is why comparison across paradigms remains difficult.
% (§\ref{subsec:p3_s8_2_summary}).


\subsubsection{Challenges and Trends}
\label{subsubsec:p3_s8_2_2_challenges}

Scene synthesis meets the problems of object synthesis on a structured set: more conditioning variables, more output components, and longer trajectories enter the cost at once.

\paragraph{Cross-Paradigm Inconsistency.}
The four paradigms solve subtly different problems, and each leaves a gap when used alone: a layout-conditioned generator can satisfy the symbolic prior $\mathbf{G}_{\text{in}}$ yet produce a non-traversable static scene, and a view-to-3D generator can produce a navigable Gaussian field that ignores the requested object semantics.
The gaps surface as \emph{layout-adherence drift}, where the generated NeRF or 3DGS no longer satisfies the input layout once the user roams outside the trained view envelope; \emph{object drop}, where small or thin objects from $\mathbf{L}$ are silently omitted during continuous-field synthesis; and \emph{geometry flattening}, where a multi-object composition collapses into a single textured slab when the 2D or video prior dominates.
Hybrids such as GALA3D~\cite{zhou2024gala3d} and ViewCrafter~\cite{yu2024viewcrafter} chain a layout or geometry prior with a Gaussian or video renderer to keep both halves honest, and future systems will more likely compose the four paradigms than choose among them.

\paragraph{Scaling Beyond Indoor Rooms.}
Almost all paradigms remain anchored to furnished indoor rooms because that is where the training corpora live: 3D-FRONT~\cite{fu20213dfront} alone supplies the layout supervision for the dominant P1 methods, and ScanNet, Matterport3D, and Hypersim dominate P2 / P3.
The result is an outdoor and urban scale gap that CityCraft~\cite{deng2024citycraft} is among the few systems to address, a category imbalance with many furniture classes but few terrain, vegetation, or industrial ones, and an open-world domain shift under which methods trained on residential interiors fail on industrial sites, excavations, or invented environments.
Web-scraped open-vocabulary scene corpora and hybrid synthetic-real pipelines that blend procedurally generated environments with captured scans are the active responses, both for breadth and to control the long tail of the scene-class distribution.

\paragraph{Camera-Trajectory Coherence under Long-Horizon Synthesis.}
The video paradigm (§\ref{para:p3_s8_2_2_video}) excels at short camera moves but degrades on long, unbounded trajectories that require persistent world memory.
Long trajectories expose \emph{temporal blur}, where textures shimmer or drift between frames; \emph{scene amnesia}, where objects reappear with different geometry when the camera revisits a location; and \emph{pose-conditioning collapse}, where the model learns to ignore the camera signal in favor of the marginal video distribution.
ViewCrafter~\cite{yu2024viewcrafter}'s unprojection and reprojection scaffold, Gen3C~\cite{ren2025gen3c}'s 3D cache, and CAT4D~\cite{wu2025cat4d}'s joint multi-view and video training all inject an explicit 3D inductive bias back into the video prior, but the trade-off between the video prior's perceptual quality and the geometric prior's rigidity remains an open design tension.

\paragraph{Evaluation.}
A P1 layout-adherence score and a P3 PSNR cannot be ranked against each other, no single dataset evaluates layout faithfulness, geometric fidelity, photometric quality, and pose adherence jointly, and no human-aligned evaluator analogous to HyperScore~\cite{zhang2025mate3d} has been validated on scene outputs.
% ; the cross-subfield picture is drawn in §\ref{subsec:p3_s8_2_summary}.

\paragraph{Trends: Toward 4D and World Models.}
Scene synthesis targets static geometry conditioned on camera poses, but the same formulation is being extended toward 4D synthesis and action-conditioned world modeling: Dreamer~\cite{hafner2019dreamer}, GameNGen~\cite{valevski2024diffusion}, Diamond~\cite{alonso2024diffusionworldmodelingvisual}, RTFM~\cite{worldlabs_rtfm}, and Genie3~\cite{genie3} move from static view or video synthesis toward open-ended environment generation in which spatial structure, temporal evolution, and interaction are modeled together.
From the standpoint of this subsection, scalable, temporally coherent video generation with explicit viewpoint control is the foundation those systems build on; they are treated as world models in §\ref{sec:p3_s8_6}.
